%%%% Начало оформления заголовка - оставить без изменений !!! %%%%
\input{my_folder/task_settings}	% настройки - начало


{\centering%

	\Ministry\\
	\SPbPU\\
	\institute
\par}
}\intervalS% завершает input

				\noindent
				\begin{minipage}{\linewidth}
				\vspace{\mfloatsep} % интервал
				\begin{tabularx}{\linewidth}{Xl}
					&УТВЕРЖДАЮ      \\
					&\HeadTitle     \\
					&\underline{\hspace*{0.1\textheight}} \Head     \\
					&«\underline{\hspace*{0.05\textheight}}» \underline{\hspace*{0.1\textheight}} \thesisYear~г.  \\
				\end{tabularx}
				\vspace{\mfloatsep} % интервал
				\end{minipage}

\intervalS{\centering\bfseries%

				ЗАДАНИЕ\\
				на выполнение %с 2020 года
				%по выполнению % до 2020 года
				выпускной квалификационной работы


\intervalS\normalfont%

				студенту \uline{\AuthorFullDat{}, гр.~\group}


\par}\intervalS%
%%%%
%%%% Конец оформления заголовка  %%%%



\begin{enumerate}[1.]
	\item Тема работы: {\expandafter \ulined <<Проектирование и реализация iOS"=при\-ложения как комплексного помощника для мусульман>>.}
	%\item Тема работы (на английском языке): \uline{\thesisTitleEn.} % вероятно после 2021 года
	\item Срок сдачи студентом законченной работы: \uline{\thesisDeadline.}
	\item Исходные данные по работе:
	\begin{enumerate}[label=\theenumi\arabic*.]
		\sloppy
		\item Коран: переводы смыслов на русский язык [Электронный ресурс] // Quran.com: [официальный сайт]. "--- URL: \url{https://quran.com/ru} (дата обращения: 01.02.2026). "--- Режим доступа: свободный.
		\item Apple Developer Documentation. Swift Programming Language [Электронный ресурс] // Apple Developer: [официальный сайт]. "--- URL: \url{https://developer.apple.com/documentation/swift}
		\item Apple Developer Documentation. SwiftUI Framework [Электронный ресурс] // Apple Developer: [официальный сайт]. "--- URL: \url{https://developer.apple.com/documentation/swiftui}
		\item PyTorch Documentation [Электронный ресурс] // PyTorch: [официальный сайт]. "--- URL: \url{https://pytorch.org/docs/stable/index.html}
		\item Hugging Face Transformers Documentation [Электронный ресурс] // Hugging Face: [официальный сайт]. "--- URL: \url{https://huggingface.co/docs/transformers}
		\item FastAPI Documentation [Электронный ресурс] // FastAPI: [официальный сайт]. "--- URL: \url{https://fastapi.tiangolo.com/}
		\item Spring Framework Documentation [Электронный ресурс] // Spring: [официальный сайт]. "--- URL: \url{https://docs.spring.io/spring-framework/docs/current/reference/html/}
	\end{enumerate}
	%\printbibliographyTask % печать списка источников % КОММЕНТИРУЕМ ЕСЛИ НЕ ИСПОЛЬЗУЕТСЯ
	\item Содержание работы (перечень подлежащих разработке вопросов):
	\begin{enumerate}[label=\theenumi\arabic*.]
		\item Анализ предметной области и обзор существующих решений для мусульманской аудитории.
		\item Формулирование функциональных и нефункциональных требований к приложению.
		\item Проектирование архитектуры системы (iOS-клиент, backend, RAG-пайплайн, БД).
		\item Разработка и интеграция модулей приложения.
		\item Тестирование работоспособности системы.
	\end{enumerate}
	\item Перечень графического материала (с указанием обязательных чертежей): \uline{Нет.}
	\item Перечень используемых информационных технологий, в том числе программное обеспечение, облачные сервисы, базы данных и прочие сквозные цифровые технологии (при наличии):
	\begin{enumerate}[label=\theenumi\arabic*.]
		\item СУБД PostgreSQL.
		\item Облачный API-сервис OpenRouter (LLM-генерация).
		\item Фреймворк FastAPI (Python).
		\item Фреймворк Spring Boot (Java).
	\end{enumerate}
	\item Консультанты по работе:
	\begin{enumerate}[label=\theenumi\arabic*.]
		\item \uline{\emakefirstuc{\ConsultantNormDegree}, \ConsultantNorm{} (нормоконтроль).}
	\end{enumerate}
	\item Дата выдачи задания: \uline{22.12.2025.}
\end{enumerate}

\intervalS%можно удалить пробел

Руководитель ВКР \uline{\hspace*{0.1\textheight} \Supervisor}


\intervalS%можно удалить пробел

%Консультант по нормоконтролю \uline{\hspace*{0.1\textheight} \ConsultantNorm}%ПОКА НЕ ТРЕБУЕТСЯ, Т.К. ОН У ВСЕХ ПО УМОЛЧАНИЮ

Задание принял к исполнению \uline{22.12.2025}

\intervalS%можно удалить пробел

Студент \uline{\hspace*{0.1\textheight}  \Author}



\input{my_folder/task_settings_restore}	% настройки - конец
